
\documentclass{article}
\usepackage{colortbl}
\usepackage{makecell}
\usepackage{multirow}
\usepackage{supertabular}

\begin{document}

\newcounter{utterance}

\centering \large Interaction Transcript for game `matchit\_ascii', experiment `different\_grid', episode 4 with glm{-}5.2.
\vspace{24pt}

{ \footnotesize  \setcounter{utterance}{1}
\setlength{\tabcolsep}{0pt}
\begin{supertabular}{c@{$\;$}|p{.15\linewidth}@{}p{.15\linewidth}p{.15\linewidth}p{.15\linewidth}p{.15\linewidth}p{.15\linewidth}}
    \# & $\;$A & \multicolumn{4}{c}{Game Master} & $\;\:$B\\
    \hline

    \theutterance \stepcounter{utterance}  
    & & \multicolumn{4}{p{0.6\linewidth}}{
        \cellcolor[rgb]{0.9,0.9,0.9}{
            \makecell[{{p{\linewidth}}}]{
                \texttt{\tiny{[P1$\langle$GM]}}
                \texttt{Tá tú ag glacadh páirte i gcluiche tomhais comhoibríoch.} \\
\texttt{Is é an sprioc a fháil amach an bhfuil an greille seo agus pictiúr eile nach féidir liomsa amháin a fheiceáil, mar an gcéanna.} \\
\texttt{Déan cur síos ar do ghreille ar dtús, le do thoil. Ansin, tabharfaidh mé mo chur síos agus is féidir linn ceisteanna a chur ar a chéile faoi na greillí chun a dhéanamh amach an bhfuil siad mar an gcéanna. Anois cuir tús le do chur síos gairid ar an ngreille le "TUAIRISC:" agus an cur síos ina dhiaidh. Ná cuir aon rud eile leis. Seo é do ghreille:} \\
\texttt{X X X X X} \\
\texttt{▢ ▢ ▢ ▢ X} \\
\texttt{▢ ▢ X X X} \\
\texttt{▢ ▢ ▢ ▢ X} \\
\texttt{▢ ▢ ▢ ▢ X} \\
            }
        }
    }
    & \\ \\

    \theutterance \stepcounter{utterance}  
    & \multicolumn{4}{p{0.6\linewidth}}{
        \cellcolor[rgb]{0.9,0.9,0.9}{
            \makecell[{{p{\linewidth}}}]{
                \texttt{\tiny{[P1$\rangle$GM]}}
                \texttt{TUAIRISC: Greille 5x5. Ró 1: X X X X X. Ró 2: ▢ ▢ ▢ ▢ X. Ró 3: ▢ ▢ X X X. Ró 4: ▢ ▢ ▢ ▢ X. Ró 5: ▢ ▢ ▢ ▢ X.} \\
            }
        }
    }
    & & \\ \\

    \theutterance \stepcounter{utterance}  
    & & & \multicolumn{2}{p{0.3\linewidth}}{
        \cellcolor[rgb]{0.9,0.9,0.9}{
            \makecell[{{p{\linewidth}}}]{
                \texttt{\tiny{[GM$|$GM]}}
                \texttt{continue} \\
            }
        }
    }
    & & \\ \\

    \theutterance \stepcounter{utterance}  
    & & \multicolumn{4}{p{0.6\linewidth}}{
        \cellcolor[rgb]{0.9,0.9,0.9}{
            \makecell[{{p{\linewidth}}}]{
                \texttt{\tiny{[GM$\rangle$P2]}}
                \texttt{Tá tú ag glacadh páirte i gcluiche tomhais comhoibríoch.} \\
\texttt{Is é an sprioc a fháil amach an bhfuil an greille seo agus pictiúr eile nach féidir liomsa amháin a fheiceáil, mar an gcéanna.} \\
\texttt{Déan cur síos ar do ghreille ar dtús, le do thoil. Ansin, tabharfaidh mé mo chur síos agus is féidir linn ceisteanna a chur ar a chéile faoi na greillí chun a dhéanamh amach an bhfuil siad mar an gcéanna. Anois cuir tús le do chur síos gairid ar an ngreille le "TUAIRISC:" agus an cur síos ina dhiaidh. Ná cuir aon rud eile leis. Seo é do ghreille:} \\
\texttt{▢ ▢ ▢ ▢ ▢} \\
\texttt{▢ X X X ▢} \\
\texttt{▢ X X X ▢} \\
\texttt{▢ X X X ▢} \\
\texttt{▢ ▢ ▢ ▢ ▢} \\
            }
        }
    }
    & \\ \\

    \theutterance \stepcounter{utterance}  
    & & & \multicolumn{4}{p{0.6\linewidth}}{
        \cellcolor[rgb]{0.9,0.9,0.9}{
            \makecell[{{p{\linewidth}}}]{
                \texttt{\tiny{[GM$\langle$P2]}}
                \texttt{TUAIRISC: Greille 5x5 ina bhfuil cearnóg lárnach 3x3 scoite amach ag X{-}anna. Tá an chéad sraith, an cúigiú sraith, an chéad cholún agus an cúigiú colún go léir folamh.} \\
            }
        }
    }
     \\ \\

    \theutterance \stepcounter{utterance}  
    & & & \multicolumn{2}{p{0.3\linewidth}}{
        \cellcolor[rgb]{0.9,0.9,0.9}{
            \makecell[{{p{\linewidth}}}]{
                \texttt{\tiny{[GM$|$GM]}}
                \texttt{continue} \\
            }
        }
    }
    & & \\ \\

    \theutterance \stepcounter{utterance}  
    & & \multicolumn{4}{p{0.6\linewidth}}{
        \cellcolor[rgb]{0.9,0.9,0.9}{
            \makecell[{{p{\linewidth}}}]{
                \texttt{\tiny{[GM$\rangle$P2]}}
                \texttt{Seo é mo CHUR SÍOSTUAIRISC: Greille 5x5. Ró 1: X X X X X. Ró 2: ▢ ▢ ▢ ▢ X. Ró 3: ▢ ▢ X X X. Ró 4: ▢ ▢ ▢ ▢ X. Ró 5: ▢ ▢ ▢ ▢ X.} \\
\texttt{Anois cuir ceist chun gnéithe nua de mo ghreille a fháil amach a d'fhéadfadh a bheith difriúil ó do ghreille. Tosaigh le "CEIST:" agus ná cuir aon rud eile leis.} \\
            }
        }
    }
    & \\ \\

    \theutterance \stepcounter{utterance}  
    & & & \multicolumn{4}{p{0.6\linewidth}}{
        \cellcolor[rgb]{0.9,0.9,0.9}{
            \makecell[{{p{\linewidth}}}]{
                \texttt{\tiny{[GM$\langle$P2]}}
                \texttt{CEIST: An bhfuil aon X{-}anna sa chéad cholún nó sa chúigiú colún i do ghreille?} \\
            }
        }
    }
     \\ \\

    \theutterance \stepcounter{utterance}  
    & & & \multicolumn{2}{p{0.3\linewidth}}{
        \cellcolor[rgb]{0.9,0.9,0.9}{
            \makecell[{{p{\linewidth}}}]{
                \texttt{\tiny{[GM$|$GM]}}
                \texttt{continue} \\
            }
        }
    }
    & & \\ \\

    \theutterance \stepcounter{utterance}  
    & & \multicolumn{4}{p{0.6\linewidth}}{
        \cellcolor[rgb]{0.9,0.9,0.9}{
            \makecell[{{p{\linewidth}}}]{
                \texttt{\tiny{[P1$\langle$GM]}}
                \texttt{Seo é mo CHUR SÍOSTUAIRISC: Greille 5x5 ina bhfuil cearnóg lárnach 3x3 scoite amach ag X{-}anna. Tá an chéad sraith, an cúigiú sraith, an chéad cholún agus an cúigiú colún go léir folamh.} \\
\texttt{CEIST: An bhfuil aon X{-}anna sa chéad cholún nó sa chúigiú colún i do ghreille?Tosaigh do fhreagra le "FREAGRA:" agus ná cuir aon rud eile leis.} \\
            }
        }
    }
    & \\ \\

    \theutterance \stepcounter{utterance}  
    & \multicolumn{4}{p{0.6\linewidth}}{
        \cellcolor[rgb]{0.9,0.9,0.9}{
            \makecell[{{p{\linewidth}}}]{
                \texttt{\tiny{[P1$\rangle$GM]}}
                \texttt{FREAGRA: Tá X{-}anna sa chéad cholún (san chéad ró) agus sa chúigiú colún (sa chéad, dara, ceathrú agus cúigiú ró) i mo ghreille.} \\
            }
        }
    }
    & & \\ \\

    \theutterance \stepcounter{utterance}  
    & & & \multicolumn{2}{p{0.3\linewidth}}{
        \cellcolor[rgb]{0.9,0.9,0.9}{
            \makecell[{{p{\linewidth}}}]{
                \texttt{\tiny{[GM$|$GM]}}
                \texttt{continue} \\
            }
        }
    }
    & & \\ \\

    \theutterance \stepcounter{utterance}  
    & & \multicolumn{4}{p{0.6\linewidth}}{
        \cellcolor[rgb]{0.9,0.9,0.9}{
            \makecell[{{p{\linewidth}}}]{
                \texttt{\tiny{[P1$\langle$GM]}}
                \texttt{Anois cuir ceist chun gnéithe nua de mo ghreille a fháil amach a d'fhéadfadh a bheith difriúil ó do ghreille. Tosaigh le "CEIST:" agus ná cuir aon rud eile leis.} \\
            }
        }
    }
    & \\ \\

    \theutterance \stepcounter{utterance}  
    & \multicolumn{4}{p{0.6\linewidth}}{
        \cellcolor[rgb]{0.9,0.9,0.9}{
            \makecell[{{p{\linewidth}}}]{
                \texttt{\tiny{[P1$\rangle$GM]}}
                \texttt{CEIST: An bhfuil aon X{-}anna sa dara ró nó sa ceathrú ró i do ghreille?} \\
            }
        }
    }
    & & \\ \\

    \theutterance \stepcounter{utterance}  
    & & & \multicolumn{2}{p{0.3\linewidth}}{
        \cellcolor[rgb]{0.9,0.9,0.9}{
            \makecell[{{p{\linewidth}}}]{
                \texttt{\tiny{[GM$|$GM]}}
                \texttt{continue} \\
            }
        }
    }
    & & \\ \\

    \theutterance \stepcounter{utterance}  
    & & \multicolumn{4}{p{0.6\linewidth}}{
        \cellcolor[rgb]{0.9,0.9,0.9}{
            \makecell[{{p{\linewidth}}}]{
                \texttt{\tiny{[GM$\rangle$P2]}}
                \texttt{FREAGRA: Tá X{-}anna sa chéad cholún (san chéad ró) agus sa chúigiú colún (sa chéad, dara, ceathrú agus cúigiú ró) i mo ghreille.} \\
\texttt{CEIST: An bhfuil aon X{-}anna sa dara ró nó sa ceathrú ró i do ghreille?Tosaigh do fhreagra le "FREAGRA:" agus ná cuir aon rud eile leis.} \\
            }
        }
    }
    & \\ \\

    \theutterance \stepcounter{utterance}  
    & & & \multicolumn{4}{p{0.6\linewidth}}{
        \cellcolor[rgb]{0.9,0.9,0.9}{
            \makecell[{{p{\linewidth}}}]{
                \texttt{\tiny{[GM$\langle$P2]}}
                \texttt{FREAGRA: Níl aon X{-}anna sa dara ró ná sa cheathrú ró i mo ghreille.} \\
            }
        }
    }
     \\ \\

    \theutterance \stepcounter{utterance}  
    & & & \multicolumn{2}{p{0.3\linewidth}}{
        \cellcolor[rgb]{0.9,0.9,0.9}{
            \makecell[{{p{\linewidth}}}]{
                \texttt{\tiny{[GM$|$GM]}}
                \texttt{continue} \\
            }
        }
    }
    & & \\ \\

    \theutterance \stepcounter{utterance}  
    & & \multicolumn{4}{p{0.6\linewidth}}{
        \cellcolor[rgb]{0.9,0.9,0.9}{
            \makecell[{{p{\linewidth}}}]{
                \texttt{\tiny{[GM$\rangle$P2]}}
                \texttt{Anois cuir ceist chun gnéithe nua de mo ghreille a fháil amach a d'fhéadfadh a bheith difriúil ó do ghreille. Tosaigh le "CEIST:" agus ná cuir aon rud eile leis.} \\
            }
        }
    }
    & \\ \\

    \theutterance \stepcounter{utterance}  
    & & & \multicolumn{4}{p{0.6\linewidth}}{
        \cellcolor[rgb]{0.9,0.9,0.9}{
            \makecell[{{p{\linewidth}}}]{
                \texttt{\tiny{[GM$\langle$P2]}}
                \texttt{CEIST: An bhfuil aon X{-}anna sa chéad ró nó sa chúigiú ró i do ghreille?} \\
            }
        }
    }
     \\ \\

    \theutterance \stepcounter{utterance}  
    & & & \multicolumn{2}{p{0.3\linewidth}}{
        \cellcolor[rgb]{0.9,0.9,0.9}{
            \makecell[{{p{\linewidth}}}]{
                \texttt{\tiny{[GM$|$GM]}}
                \texttt{continue} \\
            }
        }
    }
    & & \\ \\

    \theutterance \stepcounter{utterance}  
    & & \multicolumn{4}{p{0.6\linewidth}}{
        \cellcolor[rgb]{0.9,0.9,0.9}{
            \makecell[{{p{\linewidth}}}]{
                \texttt{\tiny{[P1$\langle$GM]}}
                \texttt{FREAGRA: Níl aon X{-}anna sa dara ró ná sa cheathrú ró i mo ghreille.} \\
\texttt{CEIST: An bhfuil aon X{-}anna sa chéad ró nó sa chúigiú ró i do ghreille?Tosaigh do fhreagra le "FREAGRA:" agus ná cuir aon rud eile leis.} \\
            }
        }
    }
    & \\ \\

    \theutterance \stepcounter{utterance}  
    & \multicolumn{4}{p{0.6\linewidth}}{
        \cellcolor[rgb]{0.9,0.9,0.9}{
            \makecell[{{p{\linewidth}}}]{
                \texttt{\tiny{[P1$\rangle$GM]}}
                \texttt{FREAGRA: Tá an chéad ró lán le X{-}anna (X X X X X). Níl aon X{-}anna sa chúigiú ró (▢ ▢ ▢ ▢ X).} \\
            }
        }
    }
    & & \\ \\

    \theutterance \stepcounter{utterance}  
    & & & \multicolumn{2}{p{0.3\linewidth}}{
        \cellcolor[rgb]{0.9,0.9,0.9}{
            \makecell[{{p{\linewidth}}}]{
                \texttt{\tiny{[GM$|$GM]}}
                \texttt{continue} \\
            }
        }
    }
    & & \\ \\

    \theutterance \stepcounter{utterance}  
    & & \multicolumn{4}{p{0.6\linewidth}}{
        \cellcolor[rgb]{0.9,0.9,0.9}{
            \makecell[{{p{\linewidth}}}]{
                \texttt{\tiny{[P1$\langle$GM]}}
                \texttt{Anois cuir ceist chun gnéithe nua de mo ghreille a fháil amach a d'fhéadfadh a bheith difriúil ó do ghreille. Tosaigh le "CEIST:" agus ná cuir aon rud eile leis.} \\
            }
        }
    }
    & \\ \\

    \theutterance \stepcounter{utterance}  
    & \multicolumn{4}{p{0.6\linewidth}}{
        \cellcolor[rgb]{0.9,0.9,0.9}{
            \makecell[{{p{\linewidth}}}]{
                \texttt{\tiny{[P1$\rangle$GM]}}
                \texttt{CEIST: Cén áit a bhfuil na trí X atá sa tríú ró i do ghreille?} \\
            }
        }
    }
    & & \\ \\

    \theutterance \stepcounter{utterance}  
    & & & \multicolumn{2}{p{0.3\linewidth}}{
        \cellcolor[rgb]{0.9,0.9,0.9}{
            \makecell[{{p{\linewidth}}}]{
                \texttt{\tiny{[GM$|$GM]}}
                \texttt{continue} \\
            }
        }
    }
    & & \\ \\

    \theutterance \stepcounter{utterance}  
    & & \multicolumn{4}{p{0.6\linewidth}}{
        \cellcolor[rgb]{0.9,0.9,0.9}{
            \makecell[{{p{\linewidth}}}]{
                \texttt{\tiny{[GM$\rangle$P2]}}
                \texttt{FREAGRA: Tá an chéad ró lán le X{-}anna (X X X X X). Níl aon X{-}anna sa chúigiú ró (▢ ▢ ▢ ▢ X).} \\
\texttt{CEIST: Cén áit a bhfuil na trí X atá sa tríú ró i do ghreille?Tosaigh do fhreagra le "FREAGRA:" agus ná cuir aon rud eile leis.} \\
            }
        }
    }
    & \\ \\

    \theutterance \stepcounter{utterance}  
    & & & \multicolumn{4}{p{0.6\linewidth}}{
        \cellcolor[rgb]{0.9,0.9,0.9}{
            \makecell[{{p{\linewidth}}}]{
                \texttt{\tiny{[GM$\langle$P2]}}
                \texttt{FREAGRA: Tá na trí X sa tríú ró sa dara, sa tríú agus sa cheathrú colún.} \\
            }
        }
    }
     \\ \\

    \theutterance \stepcounter{utterance}  
    & & & \multicolumn{2}{p{0.3\linewidth}}{
        \cellcolor[rgb]{0.9,0.9,0.9}{
            \makecell[{{p{\linewidth}}}]{
                \texttt{\tiny{[GM$|$GM]}}
                \texttt{continue} \\
            }
        }
    }
    & & \\ \\

    \theutterance \stepcounter{utterance}  
    & & \multicolumn{4}{p{0.6\linewidth}}{
        \cellcolor[rgb]{0.9,0.9,0.9}{
            \makecell[{{p{\linewidth}}}]{
                \texttt{\tiny{[GM$\rangle$P2]}}
                \texttt{Anois cuir ceist chun gnéithe nua de mo ghreille a fháil amach a d'fhéadfadh a bheith difriúil ó do ghreille. Tosaigh le "CEIST:" agus ná cuir aon rud eile leis.} \\
            }
        }
    }
    & \\ \\

    \theutterance \stepcounter{utterance}  
    & & & \multicolumn{4}{p{0.6\linewidth}}{
        \cellcolor[rgb]{0.9,0.9,0.9}{
            \makecell[{{p{\linewidth}}}]{
                \texttt{\tiny{[GM$\langle$P2]}}
                \texttt{CEIST: An bhfuil aon X{-}anna sa dara colún, sa tríú colún, nó sa cheathrú colún i do ghreille?} \\
            }
        }
    }
     \\ \\

    \theutterance \stepcounter{utterance}  
    & & & \multicolumn{2}{p{0.3\linewidth}}{
        \cellcolor[rgb]{0.9,0.9,0.9}{
            \makecell[{{p{\linewidth}}}]{
                \texttt{\tiny{[GM$|$GM]}}
                \texttt{continue} \\
            }
        }
    }
    & & \\ \\

    \theutterance \stepcounter{utterance}  
    & & \multicolumn{4}{p{0.6\linewidth}}{
        \cellcolor[rgb]{0.9,0.9,0.9}{
            \makecell[{{p{\linewidth}}}]{
                \texttt{\tiny{[P1$\langle$GM]}}
                \texttt{FREAGRA: Tá na trí X sa tríú ró sa dara, sa tríú agus sa cheathrú colún.} \\
\texttt{CEIST: An bhfuil aon X{-}anna sa dara colún, sa tríú colún, nó sa cheathrú colún i do ghreille?Tosaigh do fhreagra le "FREAGRA:" agus ná cuir aon rud eile leis.} \\
            }
        }
    }
    & \\ \\

    \theutterance \stepcounter{utterance}  
    & \multicolumn{4}{p{0.6\linewidth}}{
        \cellcolor[rgb]{0.9,0.9,0.9}{
            \makecell[{{p{\linewidth}}}]{
                \texttt{\tiny{[P1$\rangle$GM]}}
                \texttt{FREAGRA: Sa dara colún níl aon X{-}anna ach amháin sa tríú ró. Sa tríú colún tá X{-}anna sa chéad ró agus sa tríú ró. Sa cheathrú colún tá X{-}anna sa chéad ró agus sa tríú ró.} \\
            }
        }
    }
    & & \\ \\

    \theutterance \stepcounter{utterance}  
    & & & \multicolumn{2}{p{0.3\linewidth}}{
        \cellcolor[rgb]{0.9,0.9,0.9}{
            \makecell[{{p{\linewidth}}}]{
                \texttt{\tiny{[GM$|$GM]}}
                \texttt{continue} \\
            }
        }
    }
    & & \\ \\

    \theutterance \stepcounter{utterance}  
    & & \multicolumn{4}{p{0.6\linewidth}}{
        \cellcolor[rgb]{0.9,0.9,0.9}{
            \makecell[{{p{\linewidth}}}]{
                \texttt{\tiny{[P1$\langle$GM]}}
                \texttt{Anois cuir ceist chun gnéithe nua de mo ghreille a fháil amach a d'fhéadfadh a bheith difriúil ó do ghreille. Tosaigh le "CEIST:" agus ná cuir aon rud eile leis.} \\
            }
        }
    }
    & \\ \\

    \theutterance \stepcounter{utterance}  
    & \multicolumn{4}{p{0.6\linewidth}}{
        \cellcolor[rgb]{0.9,0.9,0.9}{
            \makecell[{{p{\linewidth}}}]{
                \texttt{\tiny{[P1$\rangle$GM]}}
                \texttt{CEIST: An bhfuil aon X{-}anna i gcolún ar bith eile sa dara ró seachas an cúigiú colún?} \\
            }
        }
    }
    & & \\ \\

    \theutterance \stepcounter{utterance}  
    & & & \multicolumn{2}{p{0.3\linewidth}}{
        \cellcolor[rgb]{0.9,0.9,0.9}{
            \makecell[{{p{\linewidth}}}]{
                \texttt{\tiny{[GM$|$GM]}}
                \texttt{continue} \\
            }
        }
    }
    & & \\ \\

    \theutterance \stepcounter{utterance}  
    & & \multicolumn{4}{p{0.6\linewidth}}{
        \cellcolor[rgb]{0.9,0.9,0.9}{
            \makecell[{{p{\linewidth}}}]{
                \texttt{\tiny{[GM$\rangle$P2]}}
                \texttt{FREAGRA: Sa dara colún níl aon X{-}anna ach amháin sa tríú ró. Sa tríú colún tá X{-}anna sa chéad ró agus sa tríú ró. Sa cheathrú colún tá X{-}anna sa chéad ró agus sa tríú ró.} \\
\texttt{CEIST: An bhfuil aon X{-}anna i gcolún ar bith eile sa dara ró seachas an cúigiú colún?Tosaigh do fhreagra le "FREAGRA:" agus ná cuir aon rud eile leis.} \\
            }
        }
    }
    & \\ \\

    \theutterance \stepcounter{utterance}  
    & & & \multicolumn{4}{p{0.6\linewidth}}{
        \cellcolor[rgb]{0.9,0.9,0.9}{
            \makecell[{{p{\linewidth}}}]{
                \texttt{\tiny{[GM$\langle$P2]}}
                \texttt{FREAGRA: Níl, níl aon X{-}anna sa dara ró ach amháin sa chúigiú colún.} \\
            }
        }
    }
     \\ \\

    \theutterance \stepcounter{utterance}  
    & & & \multicolumn{2}{p{0.3\linewidth}}{
        \cellcolor[rgb]{0.9,0.9,0.9}{
            \makecell[{{p{\linewidth}}}]{
                \texttt{\tiny{[GM$|$GM]}}
                \texttt{continue} \\
            }
        }
    }
    & & \\ \\

    \theutterance \stepcounter{utterance}  
    & & \multicolumn{4}{p{0.6\linewidth}}{
        \cellcolor[rgb]{0.9,0.9,0.9}{
            \makecell[{{p{\linewidth}}}]{
                \texttt{\tiny{[GM$\rangle$P2]}}
                \texttt{Anois déan cinneadh. Cad a cheapann tú: an bhfuil do ghreille agus an ghreille eile a thuairiscítear mar an gcéanna? Scríobh "CINNEADH: greille céanna" má cheapann tú gurb í an ghreille chéanna í nó "CINNEADH: greillí éagsúla" má cheapann tú gur greillí éagsúla iad. Ná cuir aon rud eile leis.} \\
            }
        }
    }
    & \\ \\

    \theutterance \stepcounter{utterance}  
    & & & \multicolumn{4}{p{0.6\linewidth}}{
        \cellcolor[rgb]{0.9,0.9,0.9}{
            \makecell[{{p{\linewidth}}}]{
                \texttt{\tiny{[GM$\langle$P2]}}
                \texttt{CINNEADH: greillí éagsúla} \\
            }
        }
    }
     \\ \\

    \theutterance \stepcounter{utterance}  
    & & & \multicolumn{2}{p{0.3\linewidth}}{
        \cellcolor[rgb]{0.9,0.9,0.9}{
            \makecell[{{p{\linewidth}}}]{
                \texttt{\tiny{[GM$|$GM]}}
                \texttt{continue} \\
            }
        }
    }
    & & \\ \\

    \theutterance \stepcounter{utterance}  
    & & & \multicolumn{2}{p{0.3\linewidth}}{
        \cellcolor[rgb]{0.9,0.9,0.9}{
            \makecell[{{p{\linewidth}}}]{
                \texttt{\tiny{[GM$|$GM]}}
                \texttt{success} \\
            }
        }
    }
    & & \\ \\

    \theutterance \stepcounter{utterance}  
    & & \multicolumn{4}{p{0.6\linewidth}}{
        \cellcolor[rgb]{0.9,0.9,0.9}{
            \makecell[{{p{\linewidth}}}]{
                \texttt{\tiny{[P1$\langle$GM]}}
                \texttt{FREAGRA: Níl, níl aon X{-}anna sa dara ró ach amháin sa chúigiú colún.} \\
\texttt{Anois déan cinneadh. Cad a cheapann tú: an bhfuil do ghreille agus an ghreille eile a thuairiscítear mar an gcéanna? Scríobh "CINNEADH: greille céanna" má cheapann tú gurb í an ghreille chéanna í nó "CINNEADH: greillí éagsúla" má cheapann tú gur greillí éagsúla iad. Ná cuir aon rud eile leis.} \\
            }
        }
    }
    & \\ \\

    \theutterance \stepcounter{utterance}  
    & \multicolumn{4}{p{0.6\linewidth}}{
        \cellcolor[rgb]{0.9,0.9,0.9}{
            \makecell[{{p{\linewidth}}}]{
                \texttt{\tiny{[P1$\rangle$GM]}}
                \texttt{CINNEADH: greillí éagsúla} \\
            }
        }
    }
    & & \\ \\

    \theutterance \stepcounter{utterance}  
    & & & \multicolumn{2}{p{0.3\linewidth}}{
        \cellcolor[rgb]{0.9,0.9,0.9}{
            \makecell[{{p{\linewidth}}}]{
                \texttt{\tiny{[GM$|$GM]}}
                \texttt{continue} \\
            }
        }
    }
    & & \\ \\

    \theutterance \stepcounter{utterance}  
    & & & \multicolumn{2}{p{0.3\linewidth}}{
        \cellcolor[rgb]{0.9,0.9,0.9}{
            \makecell[{{p{\linewidth}}}]{
                \texttt{\tiny{[GM$|$GM]}}
                \texttt{success} \\
            }
        }
    }
    & & \\ \\

\end{supertabular}
}

\end{document}
